
\medskip

%Un hôtel situé à proximité d'un site touristique dédié à la préhistoire propose deux visites dans les environs, celle d'un musée et celle d'une grotte.
%
%\medskip
%
%Une étude a montré que 70\,\% des clients de l'hôtel visitent le musée. De plus, parmi les clients visitant le musée, 60\,\% visitent la grotte.
%
%Cette étude montre aussi que 6\,\% des clients de l'hôtel ne font aucune visite.
%
%On interroge au hasard un client de l'hôtel et on note :
%
%\begin{itemize}
%\item[$\bullet~~$] $M$ l'évènement : \og le client visite le musée \fg{} ;
%\item[$\bullet~~$] $G$ l'évènement : \og le client visite la grotte \fg.
%\end{itemize}
%
%On note $\overline{M}$ l'évènement contraire de $M$,\: $\overline{G}$ l'évènement contraire de $G$, et pour tout évènement $E$, on note $p(E)$ la probabilité de $E$.
%
%Ainsi, d'après l'énoncé, on a : $p\left(\overline{M} \cap \overline{G}\right) = 0,06$.
%
%\medskip
%
%\begin{minipage}{0.53\linewidth}
\begin{enumerate}
\item 
	\begin{enumerate}
		\item %Vérifier que $p_{\overline{M}}\left(\overline{G}\right) = 0,2$, où $p_{\overline{M}}\left(\overline{G}\right)$ désigne la probabilité que le client interrogé ne visite pas la grotte sachant qu'il ne visite pas le musée.
On a $p(M) = 0,7$, donc $p\left(\overline{M}\right) = 1 - 0,7 = 0,3$.

Or $p\left(\overline{M} \cap \overline{G}\right) = 0,06$ donc 

$p\left(\overline{M} \right) \times p_{\overline{M}}\left(\overline{G}\right)=0,06 \iff 
 0,3 \times p_{\overline{M}}\left(\overline{G}\right)=0,06 \iff p_{\overline{M}}\left(\overline{G}\right) = 0,2$.

		\item ~%L'arbre pondéré ci-contre modélise la situation. Recopier et compléter cet arbre en indiquant sur chaque branche la probabilité associée.
On complète l'arbre pondéré:		
		
		
		\begin{center}
\pstree[treemode=R,nodesepA=0pt,nodesepB=2.5pt,treesep = 1cm,levelsep=2.5cm]{\TR{}}
{\pstree{\TR{$M$~}\taput{0,7}}
{\TR{$G$}\taput{0,6}
\TR{$\overline{G}$}\tbput{0,4}
}
\pstree{\TR{$\overline{M}$~}\tbput{0,3}}
{\TR{$G$}\taput{0,8}
\TR{$\overline{G}$}\tbput{0,2}
}}
\end{center}
		
		\item La probabilité de l'évènement \og le client visite la grotte et ne visite pas le musée \fg{} est:
		$p\left(\overline{M} \cap G\right) = p\left(\overline{M}\right)\times p_{\overline{M}}(M) = 0,3 \times 0,8 = 0,24$.

		\item %Montrer que $p(G) = 0,66$.
On a de même $p(M \cap G) = p(M) \times p_M(G) = 0,7 \times 0,6 = 0,42$.

D'après la loi des probabilités totales :

$p(G) = p(M \cap G) + p\left(\overline{M} \cap G\right) = 0,42 + 0,24 = 0,66$.
	\end{enumerate}

%\end{minipage}\hfill
%\begin{minipage}{0.43\linewidth}
%\begin{center}
%\pstree[treemode=R,nodesepA=0pt,nodesepB=2.5pt,treesep = 1cm,levelsep=2.5cm]{\TR{}}
%{\pstree{\TR{$M$~}\taput{\ldots}}
%{\TR{$G$}\taput{\ldots}
%\TR{$\overline{G}$}\tbput{\ldots}
%}
%\pstree{\TR{$\overline{M}$~}\tbput{\ldots}}
%{\TR{$G$}\taput{\ldots}
%\TR{$\overline{G}$}\tbput{\ldots}
%}}
%\end{center}
%\end{minipage}

\item Le responsable de l'hôtel affirme que parmi les clients qui visitent la grotte, plus de la moitié visitent également le musée.% Cette affirmation est-elle exacte ?

On calcule $p_G(M) = \dfrac{p(G \cap M}{p(G)} = \dfrac{p(M \cap G}{p(G)} = \dfrac{0,42}{0,66} = \dfrac{42}{66} = \dfrac{21}{33} = \dfrac{7}{11} \approx 0,64 > 0,5$. \\
L'affirmation est exacte.

\item %Les tarifs pour les visites sont les suivants:

%\begin{itemize}
%\item[$\bullet~~$] visite du musée : 12 euros;
%\item[$\bullet~~$] visite de la grotte : 5 euros.
%\end{itemize}
%
%On considère la variable aléatoire $T$ qui modélise la somme dépensée par un client de l'hôtel pour ces visites.
	\begin{enumerate}
		\item %Donner la loi de probabilité de $T$. On présentera les résultats sous la forme d'un tableau.
On a le tableau de probabilités suivant :
		
\begin{center}
\begin{tabularx}{\linewidth}{|l|*{4}{>{\centering \arraybackslash}X|}}\hline
évènement	&$M \cap G$	&$M \cap \overline{G}$	&$\overline{M} \cap G$	&$\overline{M} \cap \overline{G}$\rule[-3mm]{0mm}{9mm}\\ \hline
probabilité	&0,42		&0,28					&0,24					&0,06\\ \hline
dépense		&17			&12						&5						&0\\ \hline
\end{tabularx}
\end{center}

		\item %Calculer l'espérance mathématique de $T$.
D'après le tableau précédent :
		
$E(T) = 17 \times 0,42 + 12 \times 0,28 + 5 \times 0,24 + 0 \times 0,06 = 7,14 + 3,36 + 1,2 = 11,7$.

Ceci signifie que sur un grand nombre de visiteurs la dépense moyenne par visiteur est égale à 11,70~\euro.
		\item %Pour des questions de rentabilité, le responsable de l'hôtel estime que le montant
%moyen des recettes des visites doit être supérieur à $700$~euros par jour.

%Déterminer le nombre moyen de clients par journée permettant d'atteindre cet objectif.
Soit $x$ le nombre minimum de visiteurs , $x$ doit vérifier :

$11,7 \times x > 700 \iff x > \dfrac{700}{11,7}$. Or $\dfrac{700}{11,7} \approx 59,8$.

Il faut donc qu'il y ait au moins 60 visiteurs.
	\end{enumerate}
\item %Pour augmenter les recettes, le responsable souhaite que l'espérance de la variable aléatoire modélisant la somme dépensée par un client de l'hôtel pour ces visites passe à 15 euros, sans modifier le prix de visite du musée qui demeure à $12$ euros. 

%Quel prix faut-il fixer pour la visite de la grotte afin d'atteindre cet objectif ? (On admettra que l'augmentation du prix d'entrée de la grotte ne modifie pas la fréquentation des deux sites).
Soit $g$ le prix à payer pour visiter la grotte ; le tableau de probabilités devient :

\begin{center}
\begin{tabularx}{\linewidth}{|l|*{4}{>{\centering \arraybackslash}X|}}\hline
évènement&$M \cap G$&$M \cap \overline{G}$&$\overline{M} \cap G$&$\overline{M} \cap \overline{G}$
\rule[-3mm]{0mm}{9mm}\\ \hline
probabilité&0,42&0,28&0,24&0,06\\ \hline
dépense&12 + g&12&$g$&0\\ \hline
\end{tabularx}
\end{center}

L'espérance devient :

$E = 0,42(12 + g) + 12 \times 0,28 + 0,24 \times g + 0 \times 0,06 = 5,04 + 0,42g + 3,36 + 0,24g = 8,4 + 0,66g$.

Le responsable veut que :

$8,4 + 0,66g = 15 \iff 0,66g = 6,6 \iff g = 10$.

Le prix d'entrée à la grotte doit passer à 10 euros.
\item %On choisit au hasard $100$ clients de l'hôtel, en assimilant ce choix à un tirage avec remise.

%Quelle est la probabilité qu'au moins les trois quarts de ces clients aient visité la grotte à l'occasion de leur séjour à l'hôtel ?

%On donnera une valeur du résultat à $10^{-3}$ près.
Le nombre de visiteurs étant suffisamment grand pour que le tirage puisse être considéré avec remise, chaque visiteur a donc en moyenne une probabilité de 0,66 de visiter la grotte.

La variable aléatoire $G$ égale au nombre de visiteurs de la grotte suit donc une loi binomiale $\mathcal{B}(100~;~0,66)$.

Il faut donc trouver $p(G \geqslant 75)$.

La calculatrice donne $p(G\leqslant 74) \approx \np{0,9660}$, donc

$p(G \geqslant 75) = 1 - p(G\leqslant 74) \approx 1 - \np{0,9660}\approx 0,034$ au millième près.
\end{enumerate}

\bigskip

